\begin{table}[t]
\centering
\caption{Standardized Effect Sizes}
\label{tab:sde}
\begin{adjustbox}{max width=\textwidth}
\begin{tabular}{lcccccc}
\hline\hline
Outcome & $\hat{\beta}$ & SE & SD($Y$) & SDE & SE(SDE) & Classification \\
\hline
\emph{Panel A: Pooled} \\
  Total planted & $-183,797$ & $186,715$ & $155,108$ & $-0.029$ & $0.030$ & Small negative \\
  Corn & $39,717$ & $18,874$ & $72,142$ & $0.014$ & $0.006$ & Small positive \\
  Soybeans & $29,319$ & $29,854$ & $57,262$ & $0.013$ & $0.013$ & Small positive \\
  Wheat & $-252,833$ & $176,006$ & $106,301$ & $-0.058$ & $0.041$ & Moderate negative \\
\hline
\emph{Panel B: Heterogeneous (Corn)} \\
  Crop belt states & $46,884$ & $35,579$ & $94,072$ & $0.012$ & $0.009$ & Small positive \\
  Non-belt states & $35,041$ & $21,681$ & $45,931$ & $0.019$ & $0.012$ & Small positive \\
\hline\hline
\end{tabular}
\end{adjustbox}
\begin{tablenotes}[flushleft]\small
\item \textit{Notes:} \textbf{Country:} United States. \textbf{Research question:} Does the mandatory reduction of CRP acreage caps in the 2014 Farm Bill cause county-level conversion of conservation grassland to crop production? \textbf{Policy mechanism:} The Agricultural Act of 2014 mandated a step-down of the national CRP enrollment cap from 32 million to 24 million acres, forcing approximately 7 million acres of contracts to expire without reenrollment opportunity between 2013 and 2018, with county-level variation in exposure driven by pre-reform enrollment shares. \textbf{Outcome definition:} County-level planted acres for corn, soybeans, and wheat from USDA NASS QuickStats annual surveys, measuring active crop production on agricultural land. \textbf{Treatment:} Continuous; CRP acres lost between 2012--2013 and 2018--2019 averages divided by Census of Agriculture 2012 total cropland (units: share of cropland). \textbf{Data:} USDA FSA CRP county enrollment (1986--2024), USDA NASS crop acreage (2006--2022), Census of Agriculture (2012); county-year panel, 37,460 observations across 2,476 counties. \textbf{Method:} Continuous-treatment difference-in-differences with county and state-by-year fixed effects; standard errors clustered at state level (42 clusters). \textbf{Sample:} US counties with nonzero CRP enrollment and crop acreage data; excludes counties without Census 2012 cropland denominator. SDE $= \hat{\beta} \times \text{SD}(X) / \text{SD}(Y)$ where SD($X$) is the cross-sectional standard deviation of treatment intensity and SD($Y$) is the pre-treatment standard deviation of the outcome. Classification refers to magnitude, not statistical significance: Large ($|$SDE$| > 0.15$), Moderate ($0.05$--$0.15$), Small ($0.005$--$0.05$), Null ($< 0.005$).
\end{tablenotes}
\end{table}

